\documentclass[12pt,a4paper]{exam} % Adds 'answers' to print solutions
\usepackage[utf8]{inputenc}
\usepackage{amsmath}
\usepackage{amsfonts}
\usepackage{amssymb}
\usepackage{graphicx}
\graphicspath{ {figures/} }
\usepackage{todonotes}
\usepackage{hyperref}
\usepackage{xfrac}

\author{Marc Bourqui}
\title{Distributed Information Systems 2015 quizzes}

\newcommand{\mat}[1]{\ensuremath{\textbf{#1}}}

\begin{document}

\section{An Overview (week 1)}
\todo[inline]{Overview DIS}

\section{Semistructured Data Management}

\subsection{Horizontal Fragmentation}
\subsubsection{Relational Databases}
\begin{questions}

%9
\question At which phase of the database lifecycle is fragmentation performed ?
\begin{checkboxes}
\CorrectChoice At database design time
\choice During distributed query processing
\choice During updates to a distributed database
\end{checkboxes}

%10
\question The reconstruction property expresses that
\begin{checkboxes}
\choice In case of a node failure the data can be recovered from a fragment from another node
\CorrectChoice The original data can be fully recovered from the fragments
\choice Every data value of the original data can be found in at least one fragment
\end{checkboxes}
\end{questions}


\subsubsection{Primary Horizontal Fragmentation (week 2)} %Week2

\begin{questions}
%16
\question Example: application A1 accesses
\begin{enumerate}
\item Fragment F1: with frequency 3
\item Fragment F2: with frequency 1
\end{enumerate}
A1 accesses the whole relation with frequency
\missingfigure{graph}
\begin{checkboxes}
\CorrectChoice $\sfrac{13}{7}$
\choice $\sfrac{4}{7}$
\choice $\sfrac{14}{7}$
\end{checkboxes}

%26
\question Consider the access frequencies below:
How many horizontal fragments would a minimal and complete fragmentation have?
\missingfigure{table}
\begin{checkboxes}
\CorrectChoice 3
\choice 4
\choice 6
\end{checkboxes}

%27
\question Which of the following sets of simple predicates is complete?
\missingfigure{table}
\begin{checkboxes}
\choice {Location = ”Munich", Budget \textgreater\ 200000 }
\choice {Location = ”Munich", Location = ”Bangalore”}
\choice {Location = ”Paris", Budget $\le$ 200000 }
\CorrectChoice None of those
\end{checkboxes}

%33
\question Which is true for MinFrag algorithm?
\begin{checkboxes}
\choice The output is independent of the order of the input
\choice It produces a monotonically increasing set of predicates
\CorrectChoice It always terminates
\choice All of the above statements are true
\end{checkboxes}


%39
\question When deriving a horizontal fragmentation for relation $S$ from a horizontally fragmented relation $R$
\begin{checkboxes}
\CorrectChoice Some primary key attribute in $R$ must be a foreign key in $S$
\choice Some primary key attribute in $S$ must be a foreign key in $R$
\choice Both are required
\end{checkboxes}

\end{questions}

\subsection{Graph Databases (week 3)} %Week3
\subsubsection{Semi-structured Data}
\begin{questions}
%14
\question Semi-structured data
\begin{checkboxes}
\choice Is always schema-less
\CorrectChoice Always embeds schema information into the data
\choice Must always be hierarchically structured
\choice Can never be indexed
\end{checkboxes}

%15
\question Why is XML a document model?
\begin{checkboxes}
\choice It supports application-specific markup
\choice It supports domain-specific schemas
\CorrectChoice It has a serialized representation
\choice It uses HTML tags
\end{checkboxes}
\end{questions}

\subsubsection{Graph Data Model}
\begin{questions}
%22
\question In a graph database
\begin{checkboxes}
\choice There is a unique root node
\CorrectChoice Each node has a unique identifier
\choice Data values in leaf nodes are unique
\choice The labels of edges leaving a node are different
\choice There is a unique path from the root to each leaf
\end{checkboxes}

%36
\question The simulation relationship is a relation
\begin{checkboxes}
\CorrectChoice Among nodes in the data and schema graph
\choice Among edges in the data and schema graph
\choice Among sets of nodes in the data and schema graph
\choice Among sets of edges in the data and schema graph
\end{checkboxes}

%37
\question Which is true?
\begin{checkboxes}
\choice For each labelled edge in $S$ a corresponding edge in $D$ can be identified
\choice For each root node in $S$ a corresponding root node $D$ can be identified
\CorrectChoice For each leaf node in $D$ a corresponding typed node in $S$ can be identified
\choice For each node in $S$ a unique path reaching it from a root node can be identified
\end{checkboxes}

%43
\question If there exists a uniquely defined simulation relationship among a graph database $D$ and a schema graph $S$
\begin{checkboxes}
\choice The data and schema graph are simulation equivalent
\CorrectChoice Ambiguous classification cannot occur
\choice Multiple classification cannot occur
\end{checkboxes}

%48
\question If schema graph $S_1$ subsumes $S_2$
\begin{checkboxes}
\choice Every graph database corresponding to $S_1$ corresponds also to $S_2$
\CorrectChoice $S_2$ simulates $S_1$
\choice $S_1$ has fewer nodes than $S_2$
\end{checkboxes}
\end{questions}

\subsubsection{Schema Extraction}
\begin{questions}

%57
\question Which is wrong? In a dataguide
\begin{checkboxes}
\choice Every path in the data graph occurs only once
\CorrectChoice Every node in the data graph occurs only in one data guide node
\choice Every data guide node has a unique set of nodes
\choice A leaf node in the data graph corresponds always to a leaf node in the data guide
\end{checkboxes}

%63
\question In a non-deterministic schema graph
\begin{checkboxes}
\CorrectChoice Every node of the data graph occurs exactly once
\choice Every path of the data graph occurs at most once
\choice Every label of an outgoing edge of a node in the schema graph is unique
\end{checkboxes}
\end{questions}

\section{Information Retrieval and Data Mining}
\subsection{Information Retrieval (week 4)} %Week4
\subsubsection{Information Retrieval}
\begin{questions}
%12
\question A retrieval model attempts to model
\begin{checkboxes}
\choice The interface by which a user is accessing information
\CorrectChoice The importance a user gives to a piece of information
\choice The formal correctness of a query formulation by user
\choice All of the above
\end{checkboxes}

%16
\question If the top 100 documents contain 50 relevant documents
\begin{checkboxes}
\choice The precision of the system at 50 is 0.5
\CorrectChoice The precision of the system at 100 is 0.5
\choice The recall of the system is 0.5
\choice None of the above
\end{checkboxes}

%17
\question If retrieval system A has a higher precision than system B
\begin{checkboxes}
\choice The top k documents of A will have higher similarity values than the top k documents of B
\CorrectChoice The top k documents of A will contain more relevant documents than the top k documents of B
\choice A will recall more documents above a given similarity threshold than B
\choice Relevant documents in A will have higher similarity values than in B
\end{checkboxes}
\end{questions}

\subsubsection{Text-based Information Retrieval}
\begin{questions}
%25
\question Full-text retrieval means that
\begin{checkboxes}
\choice The document text is grammatically deeply analyzed for indexing
\choice The complete vocabulary of a language is used to extract index terms
\CorrectChoice All words of a text are considered as potential index terms
\choice All grammatical variations of a word are indexed
\end{checkboxes}

%26
\question The term-document matrix indicates
\begin{checkboxes}
\CorrectChoice How many relevant terms a document contains
\choice How relevant a term is for a given document
\CorrectChoice How often a relevant term occurs in a document collection
\CorrectChoice Which relevant terms are occurring in a document collection
\end{checkboxes}

%31
\question Let the query be represented by the following vectors:
(1, 0, -1) (0, -1, 1); the document by the vector (1, 0, 1)
\begin{checkboxes}
\choice Matches the query because it matches the first query vector
\CorrectChoice Matches the query because it matches the second query vector
\choice Does not match the query because it does not match the first query vector
\choice Does not match the query because it does not match the second query vector
\end{checkboxes}

%38
\question Which is right? The term frequency is normalized
\begin{checkboxes}
\CorrectChoice By the maximal frequency of a term in the document
\choice By the maximal frequency of a term in the document collection
\choice By the maximal frequency of a term in the vocabulary
\choice By the maximal term frequency of any document in the collection
\end{checkboxes}


%44
\question The inverse document frequency of a term can increase
\begin{checkboxes}
\choice By adding the term to a document that contains the term
\CorrectChoice By adding a document to a document collection that does not
contain the term
\choice By removing a document from the document collection that
does not contain the term
\choice By adding a document to a document collection that contains
the term
\end{checkboxes}

\end{questions}

\subsection{Advanced Retrieval Models (week 5)} % Week 5

\subsubsection{Latent Semantic Indexing}

\begin{questions}

% 13 1
\question In vector space retrieval each row of the matrix $M^t$ corresponds to
\begin{checkboxes}
\CorrectChoice A document
\choice A concept
\choice A query
\choice A query result
\end{checkboxes}

% 22 2
\question Applying SVD to a term-document matrix \mat{M}. Each concept is represented
\begin{checkboxes}
\choice As a singular value
\CorrectChoice As a linear combination of terms of the vocabulary
\choice As a linear combination of documents in the document collection
\choice As a least square approximation of the matrix \mat{M}
\end{checkboxes}

% 23 2
\question The number of term vectors in the SVD for LSI
\begin{checkboxes}
\choice Is smaller than the number of rows in the matrix \mat{M}
\CorrectChoice Is the same as the number of rows in the matrix \mat{M}
\choice Is larger than the number of rows in the matrix \mat{M}
\end{checkboxes}

% 32 1
\question A query transformed into the concept space for LSI has
\begin{checkboxes}
\CorrectChoice $s$ components (number of singular values)
\choice $m$ components (size of vocabulary)
\choice $n$ components (number of documents)
\end{checkboxes}
\end{questions}

\subsubsection{User Relevance Feedback}
% 40 3
\begin{questions}
\question Can documents which do not contain any keywords of the original query receive a positive similarity coefficient after relevance feedback ?
\begin{checkboxes}
\choice No
\choice Yes, independent of the values $\beta$ and $\gamma$
\CorrectChoice Yes, but only if $\beta>0$
\choice Yes, but only if $\gamma>0$
\end{checkboxes}


\subsubsection{Link-based Ranking}
% 49 1,3
\question A positive random jump value for exactly one node implies that
\begin{checkboxes}
\CorrectChoice a random walker can leave the node even without outgoing edges
\choice a random walker can reach the node multiple times even without outgoing edges
\CorrectChoice a random walker can reach the node even without incoming edges
\choice none of the above
\end{checkboxes}

% 56 3
\question Given the graph below and an initial hub vector of $(1,1,1)$. The hub-authority ranking will result in the following
\missingfigure{graph}
\begin{checkboxes}
\choice authority vector $(0,0,1$)~; hub vector $(1,1,0)$
\choice authority vector $(0,0,2)$~; hub vector $(2,2,0)$
\CorrectChoice authority vector $(0,0,1)$~; hub vector $(\frac{1}{2},\frac{1}{2},0)$
\choice authority vector $(0,0,2)$~; hub vector $(1,1,0)$
\end{checkboxes}
\end{questions}


\subsubsection{Inverted Files (week 6)}
% 10 3
\begin{questions}
\question A posting indicates
\begin{checkboxes}
\choice The frequency of a term in the vocabulary
\choice The frequency of a term in a document
\CorrectChoice The occurrence of a term in a document
\choice The list of terms occurring in a document
\end{checkboxes}

% 11 4
\question When indexing a document collection using an inverted file, the main space requirement is implied by
\begin{checkboxes}
\choice The access structure
\choice The vocabulary
\choice The index file
\CorrectChoice The postings file
\end{checkboxes}

% 18 4
\question Using a trie in index construction
\begin{checkboxes}
\choice Helps to quickly find words that have been seen before
\choice Helps to quickly decide whether a word has not seen before
\choice Helps to maintain the lexicographic order of words seen in the documents
\CorrectChoice All of the above
\end{checkboxes}

% 25 1
\question Maintaining the order of document identifiers when partitioning the document collection is important
\begin{checkboxes}
\CorrectChoice In the index merging approach for single node machines
\choice In the map-reduce approach for parallel clusters
\choice In both
\choice In neither of the two
\end{checkboxes}

\end{questions}

\subsubsection{Distributed Retrieval}
\begin{questions}

% 36 2
\question When applying Fagin’s algorithm for a query with three different terms for finding the $k$ top documents, the algorithm will scan
\begin{checkboxes}
\choice 2 different lists
\CorrectChoice 3 different lists
\choice $k$ different lists
\choice it depends how many rounds are taken
\end{checkboxes}

% 37 2
\question Once $k$ documents have been identified that occur in all of the lists
\begin{checkboxes}
\choice These are the top-$k$ documents
\CorrectChoice The top-$k$ documents are among the documents seen so far
\choice The search has to continue in round-robin till the top-$k$ documents are identified
\choice Other documents have to be searched to complete the top-$k$ list
\end{checkboxes}
\end{questions}

\section*{Credits}
Quiz questions were taken from the lecture notes of \href{http://people.epfl.ch/karl.aberer}{Prof. Karl Aberer}.%\\

%% ---------- Footer
%\hrule
%\tiny
%Rendered \today. Written by \href{http://people.epfl.ch/marc.bourqui}{Marc Bourqui}.
%\copyright Marc Bourqui. This work is licensed under the Creative Commons Attribution-ShareAlike 3.0 Unported License.
%To view a copy of this license, visit \href{http://creativecommons.org/licenses/by-sa/3.0/}{http://creativecommons.org/licenses/by-sa/3.0/} or
%send a letter to Creative Commons, 444 Castro Street, Suite 900, Mountain View, California, 94041, USA.
%\includegraphics{by-sa.png}
%
%Source code available on~: \href{https://git.epfl.ch/repo/dis15.git}{https://git.epfl.ch/repo/dis15.git}
\end{document}